\begin{table}[H]
\centering
\small
\caption{Compra e manutenção, 17/08/2017 a 03/03/2026: publicado, réplica e dados pós-T0.}
\label{tab:A1-compra-manutencao}
\begin{tabular}{lrrrrrrrr}
\toprule
Etapa & Dias & Ret. anual (\%) & Vol. (\%) & Sharpe geom. & Sharpe & Sortino & DD máx. (\%) & Maior ret. diário (\%) \\
\midrule
0. Publicado (Tab. 7.1) & -- & 38{,}7 & 68{,}7 & 0{,}56 & -- & 1{,}11 & $-$83{,}2 & -- \\
1. Réplica: código antigo, dados pré-T0 & 3.089 & 38{,}7 & 68{,}7 & 0{,}56 & 0{,}82 & 1{,}11 & $-$83{,}2 & 23{,}0 \\
2. Código antigo, dados pós-T0 & 3.120 & 38{,}3 & 68{,}3 & 0{,}56 & 0{,}82 & 1{,}10 & $-$83{,}2 & 22{,}5 \\
3. Métricas novas (Sharpe aritmético, drawdown desde 1), dados pós-T0 & 3.120 & 38{,}3 & 68{,}3 & 0{,}56 & 0{,}82 & 1{,}10 & $-$83{,}2 & 22{,}5 \\
4. Sortino usual (desvio abaixo de zero, todos os dias) & 3.120 & 38{,}3 & 68{,}3 & 0{,}56 & 0{,}82 & 1{,}20 & $-$83{,}2 & 22{,}5 \\
\bottomrule
\end{tabular}

\par\smallskip\parbox{\linewidth}{\footnotesize Antes do T0, o buraco de março/2023 entrava como um único retorno de 32 dias (23{,}0\% em 01/04/2023). Sortino: fórmula do publicado nas etapas 0 a 3; na 4, desvio abaixo de zero sobre todos os dias.}
\end{table}
